\introduction

Современные распределённые информационные системы
функционируют в условиях сетевой неоднородности,
частичных отказов узлов и возможных инфраструктурных воздействий.
Обеспечение целостности, доступности и согласованности данных
в геораспределённой среде представляет собой
одну из ключевых задач теории и практики распределённых вычислений
\cite{Tanenbaum2017}.

Особую сложность представляет межкластерная репликация,
при которой данные передаются между удаленными кластерами.
В подобных системах сетевые задержки,
потери сообщений и отказы узлов
могут приводить к нарушению согласованности
и снижению доступности сервиса.

Цель прохождения преддипломной практики —
сформулировать математическую модель геораспределённой системы
межкластерной репликации,
определить формальные критерии
целостности, доступности и согласованности данных,
а также поставить задачу выбора параметров системы,
обеспечивающих выполнение указанных критериев.

Для достижения поставленной цели необходимо решить следующие задачи:
\begin{itemize}
\item сформулировать математическую модель системы
с учётом сетевых задержек и отказов узлов;
\item определить количественные критерии целостности,
доступности и согласованности данных;
\item поставить задачу выбора параметров,
обеспечивающих выполнение указанных критериев;
\item рассмотреть возможные пути решения поставленной задачи.
\end{itemize}
